% ECONOMICS_PROBLEM: {"id": "O33-478", "title": "Technology diffusion and intangible investment", "jel_code": "O33", "source_id": "OP-0264"}
% !TeX program = xelatex
\documentclass[11pt,a4paper]{article}
\usepackage[margin=23mm]{geometry}
\usepackage{fontspec}
\setmainfont{Latin Modern Roman}
\usepackage{amsmath,amssymb,mathtools}
\usepackage{xcolor,enumitem,booktabs,longtable,array,xurl,hyperref}
\definecolor{muted}{HTML}{53636C}
\hypersetup{colorlinks=true,urlcolor=blue,linkcolor=blue}
\setlength{\parindent}{0pt}
\setlength{\parskip}{5pt}
\newcommand{\R}{\mathbb R}
\newcommand{\E}{\mathbb E}
\newcommand{\N}{\mathbb N}
\newcommand{\ind}{\mathbf 1}
\newcommand{\Var}{\operatorname{Var}}
\newcommand{\Cov}{\operatorname{Cov}}
\newcommand{\supp}{\operatorname{supp}}
\DeclareMathOperator*{\argmin}{arg\,min}
\DeclareMathOperator*{\argmax}{arg\,max}
\newcommand{\entrymeta}[3]{{\sffamily\small\color{muted}\textbf{\texttt{#1}}\quad|\quad #2\par #3\par}}
\newcommand{\Problemref}[1]{\texttt{#1}}
\newcommand{\JELref}[2]{\texttt{#2} (JEL \texttt{#1})}
\begin{document}
\section*{Technology diffusion and intangible investment}
\textbf{Problem ID:} \texttt{O33-478}\quad \textbf{JEL:} \texttt{O33}\par
% BEGIN ECONOMICS STATEMENT
\label{problem:OP-0264}
{\sffamily\small\color{muted}Original subject group: Growth and productivity\par}
\label{definitions:problem:30007545}
\textbf{Audit classification:} Published research agenda. AI diffusion is an agenda; non-AI intangible-accounting decomposition requires separate support and is editorial here. Verification concerns the cited version, not first priority or proof that this exact editorial specification remains unsolved on October 8, 2026.
\subsection*{Definitions and precise research target}
\paragraph{Editorial mathematical specification.} Fix a technology breakthrough introduced at date zero and a population of firms. Let \(d_{it}\in[0,1]\) be adoption intensity, \(J_{it}\) investment in complementary intangible capital, and \(K^J_{i,t+1}=(1-\delta)K^J_{it}+J_{it}\) its stock, with depreciation \(\delta\in[0,1]\). True output obeys a specified production function using technology, physical inputs, and \(K^J\); measured output and measured capital may omit some declared intangible expenditures. Adoption costs, financing constraints, and organizational adjustment have separately parameterized laws.
Define the date-\(H\) true productivity loss from delayed adoption and the measured-productivity discrepancy as
\[
D_H=\log A_H^{\mathrm{instant}}-\log A_H^{\mathrm{actual}},\qquad
M_H=\log A_H^{\mathrm{actual}}-\log \widetilde A_H^{\mathrm{actual}},
\]
where instantaneous adoption is feasible only under stated complementary-input and resource constraints, and \(\widetilde A\) is productivity calculated with the observed accounting convention. Identify these two quantities and their interactions from adoption histories, intangible-investment measurement, and credible variations in implementation costs. Do not count the same omitted investment simultaneously as both a diffusion barrier and a measurement error without an explicit allocation rule. Test the decomposition against later measured productivity realization. The target is technology- and horizon-specific attribution of productivity lags, not a generic assertion that intangible investment explains every productivity slowdown.

True productivity means a residual using a declared full input inventory, whereas measured productivity uses the stated published accounting convention; both require positive outputs and inputs. Construction of complementary intangible capital uses real resources, so its formation belongs either in investment output or a disclosed expensing convention. Define the instantaneous-adoption comparison through a feasible accelerated deployment policy with a stated resource cost, not free immediate installation. \(D_H\) need not be nonnegative: accelerated adoption can divert resources enough to lower finite-horizon output or change measured TFP. The accounting discrepancy \(M_H\) is separate from the difference between GDP levels; omitting both investment output and a capital input can affect numerator and denominator in opposite directions. The general technology/intangible extension is editorial; the cited agenda directly motivates AI diffusion.
\subsection*{Variants and assumption changes}
The following are \emph{editorial variants}, not additional published open conjectures.
\begin{enumerate}
\item \emph{Accounting-only change.} Hold the real adoption, employment, and investment path fixed and recompute output and capital under expensing versus capitalization. Identify \(M_H\) without attributing the accounting change to faster diffusion.
\item \emph{Costly faster diffusion.} Introduce a feasible policy reducing implementation delays at an explicit fiscal/resource cost. Compare true productivity and welfare, holding the accounting convention fixed; efficient diffusion is not necessarily instantaneous diffusion.
\end{enumerate}
\subsection*{References and source audit}
\begin{itemize}
\item Erik Brynjolfsson; Anton Korinek; Ajay Agrawal. \emph{The Economics of Transformative AI: A Research Agenda}. 2024. Locator: Section3.1, Economic Growth, printed p.10, adoption/diffusion paragraph; general intangible accounting is an editorial extension. Primary text checked. \url{https://digitaleconomy.stanford.edu/app/uploads/2024/11/ETAI-White-Paper.pdf}.
\end{itemize}
AI diffusion is an agenda; non-AI intangible-accounting decomposition requires separate support and is editorial here.
\begin{enumerate}\raggedright
\item\hypertarget{definitions:source:30007545:1}{} Erik Brynjolfsson; Anton Korinek; Ajay Agrawal. 2024. \emph{The Economics of Transformative AI: A Research Agenda}. Section3.1, Economic Growth, printed p.10, adoption/diffusion paragraph; general intangible accounting is an editorial extension.
\url{https://digitaleconomy.stanford.edu/app/uploads/2024/11/ETAI-White-Paper.pdf}.
\end{enumerate}

\subsection*{Common conventions: Source mathematical conventions}

These conventions apply to each entry unless explicitly replaced.

\textbf{Models and identification.} A primitive model \(m\) specifies all probability laws, feasible choices, information, equilibrium and measurement rules used by its target. The admissible class \(\mathcal M\) is fixed independently of outcomes. If \(P_O(m)\) is its observable probability law and \(\tau(m)\) the target, its identified set at observable law \(P\) is
\[
 \mathcal I_\tau(P)=\{\tau(m):m\in\mathcal M,\ P_O(m)=P\}.
\]
Point identification means that this set is a singleton. An empty set signals incompatibility of data and assumptions. Coordinate bounds need not describe the joint set. Bounds are sharp only if every retained target is attainable or arbitrarily approachable by compatible models. Finite bounds require explicit restrictions on unobserved quantities; unrestricted confounding, hidden wealth, extrapolation or arbitrary equilibrium selection can yield uninformative sets.

\textbf{Causal interventions.} A potential outcome \(Y_i(p)\) is the outcome under a complete policy regime \(p\), including timing, financing, information and market spillovers. The target population and units are fixed before treatment. With interference, use the complete assignment vector or a justified exposure mapping; individual no-interference notation alone cannot identify an economy-wide intervention. A difference of mean potential outcomes does not require the cross-world joint distribution. Conditioning on post-treatment migration, survival, enrollment or employment changes the estimand and needs separate justification.

\textbf{Statistical statements.} An estimator, confidence set, forecast or test specifies a sampling law, model class and loss/coverage criterion. Uniform \((1-\alpha)\)-coverage means \(\inf_{m\in\mathcal M}\Pr_m\{\tau(m)\in C_n\}\geq1-\alpha\), or an explicitly asymptotic version. The whole parameter space trivially has coverage, so informativeness requires a stated width, risk or power objective. A forecasting improvement is relative to a named benchmark, information set and evaluation population.

\textbf{Equilibrium and optimization.} An equilibrium comprises feasible plans, optimality, beliefs where needed, budget constraints and all market/resource clearing. The primitives or a stated benchmark define each correspondence \(\mathcal E(p;m)\). Nonemptiness is assumed or proved, never inferred just from compact policy sets. A selected-equilibrium target uses a declared selection rule; a robust target quantifies over all permitted equilibria. Use suprema or infima unless attainment is established. An \(\varepsilon\)-optimal policy is feasible and within \(\varepsilon\) of the finite optimum value. Report an empty feasible set separately.

\textbf{Welfare and accounting.} Normative utility, interpersonal normalization, social weights, numeraire, discounting and the use of public receipts are primitives. Behavioral utility need not coincide with the welfare criterion. Household welfare includes ownership income and rents once; adding profits again to compensating variation generally double-counts them. A monetary equivalent is defined by an explicit transfer and price convention, with monotonicity and range conditions. Average effects, mechanism contributions, transfers and welfare losses are different targets. Infinite sums require convergence and dynamic feasible plans require terminal or no-Ponzi conditions.

\textbf{Source versions.} A working-paper date, revision date and journal publication year are distinguished. A DOI for a journal version is not automatically the DOI of an earlier manuscript. Locators refer to the stated version and printed pages where specified. A metadata-only check does not verify a claim about a full-text conclusion. Every entry's source subsection narrows the provenance claim accordingly.
% END ECONOMICS STATEMENT
\end{document}
